% !TEX root = ../../main.tex

\section{论文结构安排}

本文共分为七章，各章内容安排如下。

第一章为引言，主要介绍课题研究背景、研究意义、国内外研究现状、可行性分析、本文主要研究内容以及论文结构安排，明确系统的研究目标和实现边界。

第二章为相关技术分析，主要介绍系统开发和识别任务中所采用的关键技术，包括 HarmonyOS、ArkTS、ArkUI、Spring Boot、Vue、MySQL、Redis、MinIO、MediaPipe、TensorFlow、FastAPI 以及语义修正相关技术等。

第三章为系统需求分析，主要从普通用户、后台管理员和识别服务维护三个角度分析系统功能需求，并对系统的性能、安全性、可靠性、易用性和可维护性等非功能性需求进行说明。

第四章为系统设计，主要介绍 HearBridge 系统的总体架构设计、功能模块设计、数据库设计、接口设计和关键业务流程设计，重点说明移动端、服务端、后台管理端和 Python 识别服务之间的协作关系。

第五章为系统实现，主要说明系统核心功能的具体实现过程，包括移动端训练流程实现、服务端业务接口实现、后台管理端资源维护实现、手势识别服务实现、句子视频识别链路实现以及语义修正结果组织实现。

第六章为系统测试，主要对系统核心功能和识别实验进行测试与分析，包括功能测试、接口测试、识别服务测试、句子视频识别实验、非功能性测试和系统运行效果分析。

第七章为总结与展望，主要总结本文完成的系统设计与实现工作，分析系统当前存在的不足，并对后续连续手语识别、手语辅助翻译、数字人联动展示、模型优化和训练反馈优化等方向进行展望。
